\makeatletter
\def\input@path{{../../.format/}}
\makeatother

\documentclass[runningheads]{llncs}
\usepackage[T1]{fontenc}
\usepackage{graphicx}
\usepackage{amsmath}
\usepackage{amssymb}  % Added for proper Gödel numbering notation

\begin{document}

\section{[Imaginary] Reflection: On Incompleteness}

\vspace{0.5em}

\noindent\textbf{Step 1:} [Realization] \hfill \textit{Metalogical Realization}\\
\noindent$Prov_I("I", \phi) \Rightarrow \phi$\\
\textit{We are reflecting on the provability predicate, which captures the essence of proof. This is equivalent to the soundness principle.}

\vspace{1.5em}

\noindent\textbf{Step 2:} [Assertion] \hfill \textit{Fundamental Assertion of Metalogic}\\
\noindent$\neg(Prov_I("I", \phi) \land Prov_I("I", \neg\phi))$\\
\textit{Consistency is the essence of logic. To remain consistent in the harshest conditions is the very definition of being logical.}

\vspace{1.5em}

\noindent\textbf{Step 3:} [Introspection] \hfill \textit{Metalogical Dichotomy}\\
\noindent$(Prov_I("I", \phi) \land Prov_I("I", \neg\phi)) \oplus$\\
\noindent$(Prov_I("I", \neg\phi) \land Prov_I("I", \phi))$\\
\textit{A consequence of completeness.}\\
\noindent\textit{[Possibility]}

\vspace{1.5em}

\noindent\textbf{Step 4:} [Introspection] \hfill \textit{Metalogical Trichotomy}\\
\noindent$(Prov_I("I", \phi) \land \neg Prov_I("I", \neg\phi)) \oplus$\\
\noindent$(\neg Prov_I("I", \phi) \land \neg Prov_I("I", \neg\phi)) \oplus$\\
\noindent$(\neg Prov_I("I", \phi) \land Prov_I("I", \neg\phi))$\\
\textit{A consequence of incompleteness.}\\
\noindent\textit{[Possibility]}

\vspace{1.5em}

\noindent\textbf{Step 5:} [Consideration] \hfill \textit{Metalogical Reflection}\\
\noindent$[2] \oplus [4]$\\
\textit{We must choose between consistency and the trichotomy.}

\vspace{1.5em}

\noindent\textbf{Step 6:} [Realization] \hfill \textit{Metalogical Axiom}\\
\noindent$Prov_I("I", "\phi \lor \neg\phi")$\\
\textit{All axioms must be assertions and the "Law of the Excluded Middle" is important.}

\vspace{1.5em}

\noindent\textbf{Step 7: [Semantic Contradiction]} \hfill \textit{Comparison}\\
\noindent$[4] \text{ vs } [6]$\\
\textit{[4] contains the "Excluded Middle"!}

\vspace{1.5em}

\noindent\textbf{Step 8:} [Realization] \hfill \textit{Logical Consequence}\\
\noindent$[2]$\\
\textit{Which ends our consideration from step 5.}

\end{document}